\documentclass[11pt]{amsart}
\usepackage{amsmath,amssymb,amsthm}
\usepackage[margin=1.08in]{geometry}
\usepackage[hypertexnames=false]{hyperref}
\hypersetup{hidelinks}
\usepackage[final,expansion=false]{microtype}
\setlength{\parskip}{4pt plus 1pt}
\emergencystretch=2.5em

\newtheorem{theorem}{Theorem}[section]
\newtheorem{proposition}[theorem]{Proposition}
\newtheorem{lemma}[theorem]{Lemma}
\newtheorem{corollary}[theorem]{Corollary}
\numberwithin{equation}{section}

\newcommand{\C}{\mathbb C}
\newcommand{\R}{\mathbb R}
\newcommand{\Z}{\mathbb Z}
\newcommand{\Q}{\mathbb Q}
\newcommand{\A}{\mathbb A}
\newcommand{\Rel}{\operatorname{Re}}
\newcommand{\Iml}{\operatorname{Im}}
\newcommand{\Res}{\operatorname*{Res}}
\newcommand{\supp}{\operatorname{supp}}
\newcommand{\vol}{\operatorname{vol}}
\newcommand{\rd}{\operatorname{rad}}
\newcommand{\Inf}{\operatorname{Inf}}
\newcommand{\DA}{\mathcal D_A}
\newcommand{\DAr}{\mathcal D_A^{\mathrm{rad}}}
\newcommand{\DAl}{\mathcal D_A^{\lim}}
\newcommand{\CS}{C_S}
\newcommand{\cz}{\mathfrak c_0}
\newcommand{\Tl}{T_{\lim}}
\newcommand{\Il}{I_{\lim}}

\begin{document}

\title{An explicit formula for the trace on the idele class group}

\author{David Alejandro Trejo Pizzo}
\thanks{Independent Researcher. \texttt{dtrejopizzo@gmail.com}}

\subjclass[2020]{Primary 11M06; Secondary 11M26, 11R56}
\keywords{Idele class group, trace, explicit formula, Riemann zeta function,
Tate functionals}

\begin{abstract}
Let $C_S$ be the idele class group of $\Q$ for a finite set $S$ of places
containing $\infty$, and let $T_S=\sum_{v\in S}W_v$ be the trace on
$C_c^\infty(C_S)$ whose local terms are regularized integrals over the
images of the groups $\Q_v^\times$.  For the radial functions $\Gamma_h(w)=h(-\log|w|_S)$,
$h\in C_c^\infty(\R)$, and the limit $T_{\lim}$ of $T_S$ over $S$, we prove
$T_{\lim}(\Gamma_h)=\Phi_h(0)+\Phi_h(1)-\sum_\rho\Phi_h(\rho)$, where
$\Phi_h(s)=\int h(t)e^{(s-1/2)t}\,dt$ and $\rho$ runs over the zeros of the
completed zeta function $\xi$, by comparison with Weil's explicit formula,
which we prove in this normalization.  Adding a term $c\,F(1)$ to the
trace adds $c\,h(0)$ to the right-hand side; so the identity holds for all
$h$ only if $c=0$.  For even $h$, $T_{\lim}(\Gamma_h)$ is the integral of $\Phi_hA/\pi i$ over a line $\Rel s=b>1$, with
$A(s)=\tfrac12\log\pi-\tfrac12\psi(s/2)-\zeta'/\zeta(s)$, and the sum of the
residues of $\Phi_hA$; those at $s=0,1$ are the Tate functionals of
$\Gamma_h$.  If $\xi$ has a zero off the critical line, there are real $F$
with $\Phi_F(0)=\Phi_F(1)=0$ and
$T_{\lim}(\Gamma_F\star\Gamma_F^\dagger)>0$.  For each finite $S$ the form
$(F,G)\mapsto T_S(F\star G^\dagger)$ is nondegenerate on radial functions.
\end{abstract}

\maketitle

\section{Introduction}
\label{sec:intro}

Let $S=\{\infty\}\cup S_f$ be a finite set of places of $\Q$, with $S_f$ a
nonempty set of primes, and let
\[
  \CS\;\cong\;\R^\times\times\prod_{p\in S_f}\Z_p^\times ,\qquad
  |w|_S=|w_\infty| ,
\]
be the corresponding idele class group.  In \cite{P43} the convolution
algebra $\DA(S)=C_c^\infty(\CS)$, with the involution
$F^\dagger(w)=\overline{F(w^{-1})}$, is equipped with the functionals
\[
  T^{(c)}_S(F)=\sum_{v\in S}W_v(F)+c\,F(1)\qquad(c\in\R),
\]
where $W_v(F)$ is a regularization of the integral of $F$ over the image of
$\Q_v^\times$ in $\CS$ against the density $|x|_v^{1/2}/|1-x|_v$.  These
local terms are those of Weil's explicit formula \cite{Weil1952}, written
idelically as in \cite{Connes99}.  The form
$I^{(c)}_S(F,G)=T^{(c)}_S(F\star G^\dagger)$ is Hermitian and nondegenerate,
and the two Tate functionals
$\lambda_\mp(F)=\int_{\CS}F(u)|u|_S^{\mp1/2}\,d^*u$ (the Tate characters
of \cite[Theorem~5.11]{P43}) define a homomorphism
$M=(\lambda_-,\lambda_+)$ of $\DA(S)$ onto $\C^2$, which carries the
Hermitian form $\beta_T(z,w)=z_-\overline{w_+}+z_+\overline{w_-}$.  When $S$
increases, the functionals $T^{(c)}_S$ stabilize on every element, and
define a functional $T^{(c)}_{\lim}$ and a form $I^{(c)}_{\lim}$ on the
union of the algebras $\DA(S)$.  We write $T_S$, $\Tl$, $I_S$, $\Il$ when
$c=0$.  All of this is recalled, with references, in
Section~\ref{sec:carrier}.

In this paper we compute these functionals on the functions that depend only
on $|w|_S$.  For $h\in C_c^\infty(\R)$ put
\[
  \Gamma_h(w)=h(-\log|w|_S),\qquad
  \Phi_h(s)=\int_\R h(t)\,e^{(s-\frac12)t}\,dt ;
\]
the functions $\Gamma_h$ are the elements of the algebra $\DAr$ of
\cite[Definition~3.6, Lemma~3.7]{P44}; we call them radial classes.  Let
$\xi(s)=\tfrac12s(s-1)\pi^{-s/2}\Gamma(s/2)\zeta(s)$, and let $\rho$ run
over the zeros of $\xi$, repeated according to multiplicity.

\begin{theorem}[Theorem~\ref{thm:comparison} and
Corollary~\ref{cor:c-unique}]
\label{thm:introA}
For every $h\in C_c^\infty(\R)$ the series $\sum_\rho\Phi_h(\rho)$ converges
absolutely, and
\[
  \Tl(\Gamma_h)=\Phi_h(0)+\Phi_h(1)-\sum_\rho\Phi_h(\rho).
\]
The same formula holds for $T_S(\Gamma_h)$ if $h$ vanishes outside
$(-\log\ell,\log\ell)$, where $\ell$ is the least prime not in $S_f$.  For
every real $F\in C_c^\infty(\R)$ and every $c\in\R$,
\[
  I^{(c)}_{\lim}(\Gamma_F,\Gamma_F)
  =\beta_T\bigl(M(\Gamma_F),M(\Gamma_F)\bigr)
   -\sum_\rho\Phi_F(\rho)\Phi_F(1-\rho)+c\int_\R F(t)^2\,dt ,
\]
and $\beta_T(M(\Gamma_F),M(\Gamma_F))=2\,\Phi_F(0)\Phi_F(1)$.
\end{theorem}

Thus $c=0$ is the only value of the constant for which the form
$I^{(c)}_{\lim}$ on radial classes is the sum of the Tate term and the sum
over the zeros.  The proof of Theorem~\ref{thm:introA} is a comparison of
the formula for $\Tl(\Gamma_h)$ that follows from \cite{P43,P44}
(Lemma~\ref{lem:radial-trace}) with the explicit formula of Weil
\cite{Weil1952}.  We need the latter for arbitrary $h\in C_c^\infty(\R)$,
with the archimedean term in the form $\Rel\psi(\frac14+\frac{i\tau}2)-\log\pi$
and with absolute convergence of the sum over the zeros, and we prove it in
this form in Section~\ref{sec:EF} (Theorem~\ref{thm:EF}); for the formula
under weaker hypotheses on $h$ see Barner \cite{Barner1981}.

The second result expresses the same value by residues.  Put
\[
  a(s)=\tfrac12\log\pi-\tfrac12\psi(s/2),\qquad
  f_p(s)=\frac{(\log p)\,p^{-s}}{1-p^{-s}},\qquad
  A_S=a+\sum_{p\in S_f}f_p,\qquad A=a-\frac{\zeta'}{\zeta},
\]
where $\psi=\Gamma'/\Gamma$, so that $A(s)=a(s)+\sum_pf_p(s)$ for
$\Rel s>1$.  The restriction to $|w|_S^{i\tau}$ of the symbol of $T_S$, in
the sense of \cite{P43}, is $A_S(s)+A_S(1-s)$ with $s=\frac12+i\tau$
(Lemma~\ref{lem:localsymbols}).

\begin{theorem}[Theorem~\ref{thm:residues} and
Proposition~\ref{cor:noncommute}]
\label{thm:introB}
Let $h\in C_c^\infty(\R)$ be even.
\begin{enumerate}
\item[(i)] For every $b>0$, respectively every $b>1$,
\[
  T_S(\Gamma_h)=\frac1{\pi i}\int_{\Rel s=b}\Phi_h(s)A_S(s)\,ds,\qquad
  \Tl(\Gamma_h)=\frac1{\pi i}\int_{\Rel s=b}\Phi_h(s)A(s)\,ds .
\]
\item[(ii)] The function $A$ satisfies $A(1-s)=-A(s)$.  Its poles are $s=0$
and $s=1$, simple with residue $1$, and the zeros of $\xi$, simple with
residue minus the multiplicity.  The number $\Tl(\Gamma_h)$ is the sum of
the residues of $\Phi_hA$, and the residues at $s=0$ and $s=1$ are
$\Phi_h(0)=\lambda_+(\Gamma_h)$ and $\Phi_h(1)=\lambda_-(\Gamma_h)$.
\item[(iii)] For finite $S$ the function $A_S$ is holomorphic in $\Rel s>0$
and has residue $|S_f|+1$ at $s=0$.  As $S_f$ increases to the set of all
primes, $A_S\to A$ uniformly on $\Rel s\ge b$ for every $b>1$, while
$A_S(\sigma)\to+\infty$ for every $\sigma\in(0,1]$.
\end{enumerate}
\end{theorem}

The remaining two results concern the form on radial classes.  Call
$F\in C_c^\infty(\R)$ primitive if $\Phi_F(0)=\Phi_F(1)=0$, that is, if
$M(\Gamma_F)=0$.  By Theorem~\ref{thm:introA}, for real primitive $F$ one
has $\Il(\Gamma_F,\Gamma_F)=-\sum_\rho\Phi_F(\rho)\Phi_F(1-\rho)$.  Weil's
criterion \cite{Weil1952} states that all zeros of $\xi$ lie on the line
$\Rel s=\frac12$ if and only if
$\sum_\rho\Phi_F(\rho)\overline{\Phi_F(1-\bar\rho)}\ge0$ for all test
functions $F$ of a suitable class; for real $F$ this sum is
$\sum_\rho\Phi_F(\rho)\Phi_F(1-\rho)$.  The next theorem shows that in this
criterion the test functions may be restricted to real primitive functions
in $C_c^\infty(\R)$ of either parity.

\begin{theorem}[Theorem~\ref{thm:sufficiency}]
\label{thm:introC}
Suppose that $\xi$ has a zero $\rho_0$ with $\Rel\rho_0\ne\frac12$.  Then
there are nonzero real primitive functions $F_+,F_-\in C_c^\infty(\R)$, with
$F_+$ even and $F_-$ odd, such that
\[
  \sum_\rho\Phi_{F_\pm}(\rho)\Phi_{F_\pm}(1-\rho)<0,\qquad\text{that is,}\qquad
  \Il(\Gamma_{F_\pm},\Gamma_{F_\pm})>0 .
\]
Conversely, if every zero of $\xi$ lies on the line $\Rel s=\frac12$, then
$\sum_\rho\Phi_F(\rho)\Phi_F(1-\rho)\ge0$ for every real
$F\in C_c^\infty(\R)$.
\end{theorem}

The functions $F_\pm$ are of the form $q(i\,d/dt)\,\eta^{\star N}$, where
$\eta$ is an even bump function, $\eta^{\star N}$ is its $N$-fold convolution
power for a large integer $N$, and $q$ is a polynomial determined by $N$,
$\eta$, $\rho_0$ and the finitely many zeros of height below a bound
depending on $\rho_0$ and $\eta$.  The only information on the distribution
of the zeros used in the proof is the convergence of
$\sum_\rho|\rho|^{-2}$.

\begin{theorem}[Theorem~\ref{thm:radical}]
\label{thm:introD}
Let $S$ be finite and $c\in\R$.  If $h\in C_c^\infty(\R)$ and
$I^{(c)}_S(\Gamma_h,\Gamma_g)=0$ for every $g\in C_c^\infty(\R)$, then
$h=0$.
\end{theorem}

In \cite{P43} the form $I^{(c)}_S$ is shown to be nondegenerate on
$\DA(S)$; Theorem~\ref{thm:introD} says that its restriction to the
subalgebra of radial classes is nondegenerate as well.

Section~\ref{sec:carrier} recalls what is used from \cite{P43,P44}.
Section~\ref{sec:EF} contains the explicit formula,
Section~\ref{sec:comparison} the proof of Theorem~\ref{thm:introA},
Section~\ref{sec:polar} that of Theorem~\ref{thm:introB}, and
Section~\ref{sec:sufficiency} those of Theorems~\ref{thm:introC}
and~\ref{thm:introD}.

\section{Background}
\label{sec:carrier}

Throughout, $S$ is a finite set of places of $\Q$ containing $\infty$, with
$S_f:=S\setminus\{\infty\}$ nonempty; such a set is called a \emph{chart}.
Let $\A_S^\times=\R^\times\times\prod_{p\in S_f}\Q_p^\times$, let $\Gamma_S$ be
the group of positive rational numbers whose numerator and denominator are
products of primes of $S_f$, embedded diagonally in $\A_S^\times$, and let
$|x|_S=\prod_{v\in S}|x_v|_v$.  By \cite[Lemma~2.3]{P43},
\[
  \CS:=\Gamma_S\backslash\A_S^\times\;\cong\;
  \R^\times\times\prod_{p\in S_f}\Z_p^\times ,
\]
and in these coordinates $|w|_S=|w_\infty|$.  The Haar measure $d^*w$ of $\CS$
is the product of $dx/(2|x|)$ on $\R^\times$ and of the probability Haar
measures of the groups $\Z_p^\times$, so that
\begin{equation}
\label{eq:disint}
  d^*w=dt\otimes\mu,\qquad t=\log|w|_S ,
\end{equation}
with $\mu$ the probability Haar measure of the compact group
$K_S:=\{\pm1\}\times\prod_{p\in S_f}\Z_p^\times$
\cite[(8) and the lines following]{P43}.  The algebra
$\DA(S)=C_c^\infty(\CS)$ consists of the compactly supported functions on
$\CS$ that are smooth in the archimedean coordinate and locally constant in
the others \cite[Definition~2.4]{P43}.  With
\[
  (F\star G)(w)=\int_{\CS}F(u)\,G(u^{-1}w)\,d^*u,\qquad
  F^\dagger(w)=\overline{F(w^{-1})},
\]
it is a commutative $\dagger$-algebra.

For $v\in S$ let $j_v\colon\Q_v^\times\to\CS$ send $x$ to the class of the
idele with $x$ at $v$ and $1$ at the other places of $S$, and let
$\tau_v(x)=|x|_v^{1/2}/|1-x|_v$ for $x\ne1$.  The Haar measure $d^*x$ of
$\Q_p^\times$ is normalized by $\vol_{d^*x}(\Z_p^\times)=\log p$, and that of
$\R^\times$ is $dx/(2|x|)$.

\begin{proposition}[{\cite[Definition~5.4, Convention~5.5, Theorem~5.6,
Proposition~6.3, Theorem~6.5]{P43}}]
\label{prop:trace}
For each $v\in S$ there is a linear functional $W_v$ on $\DA(S)$, defined in
\cite[Definition~5.4]{P43} by regularizing the integral of
$(F\circ j_v)\,\tau_v$ over $\Q_v^\times$ at $x=1$, such that
\[
  W_v(F)=\int_{\Q_v^\times}F(j_v(x))\,\tau_v(x)\,d^*x ,
\]
the integral being absolutely convergent, whenever $F$ vanishes on a
neighbourhood of the identity class \cite[(32) and the proof of
Theorem~5.6]{P43}.  For $c\in\R$ put
\begin{equation}
\label{eq:tracec}
  T^{(c)}_S(F):=\sum_{v\in S}W_v(F)+c\,F(1),\qquad T_S:=T^{(0)}_S .
\end{equation}
For every $c$ the form $I^{(c)}_S(F,G):=T^{(c)}_S(F\star G^\dagger)$ is
sesquilinear, Hermitian and nondegenerate on $\DA(S)$.  We write
$I_S:=I^{(0)}_S$.
\end{proposition}

In \cite{P43} the constant $c$ is fixed and omitted from the notation: the
functional $T^{(c)}_S$ is denoted $T$ there, and the form $I^{(c)}_S$ is
denoted $I_A$.

The characters of $\CS$ are
$\chi(w)=\operatorname{sgn}(w_\infty)^{\varepsilon}|w_\infty|^{i\tau}
\prod_{q\in S_f}\kappa_q(w_q)$ with $\varepsilon\in\{0,1\}$, $\tau\in\R$ and
$\kappa_q$ a character of finite order of $\Z_q^\times$.  The dual group
$\widehat{\CS}$ is thus a disjoint union of lines with coordinate $\tau$, its
\emph{components}; the \emph{principal component} $\cz$ is the one with
$\varepsilon=0$ and all $\kappa_q$ trivial, and consists of the characters
$w\mapsto|w|_S^{i\tau}$.  We put
$\widehat F(\chi)=\int_{\CS}F(w)\overline{\chi(w)}\,d^*w$ and
$d\chi=d\tau/2\pi$ on each component \cite[(39), (40) and the text after (42)]{P43}.

\begin{proposition}[{\cite[Theorem~5.12 and (67)]{P43}}]
\label{prop:symbol}
There is a real-valued function
$\widehat T_S=\widehat W_\infty+\sum_{p\in S_f}\widehat W_p$ on
$\widehat{\CS}$, real-analytic on each component, with the following
properties.
\begin{enumerate}
\item[(a)] On $\cz$, writing $\widehat W_\infty(0,\tau)$ for
$\widehat W_\infty(\cz,\tau)$,
\begin{equation}
\label{eq:symbolclosed}
  \widehat W_p(\cz,\tau)=2(\log p)\,\Rel\frac{p^{-1/2+i\tau}}{1-p^{-1/2+i\tau}},
  \qquad
  \widehat W_\infty(0,\tau)=\log\pi-\Rel\psi\bigl(\tfrac14+\tfrac{i\tau}2\bigr),
\end{equation}
where $\psi=\Gamma'/\Gamma$.  Each $\widehat W_p(\cz,\cdot)$ is bounded, and
$\widehat W_\infty(0,\tau)=-\log|\tau|+O(1)$ as $|\tau|\to\infty$.
\item[(b)] For every $c\in\R$ the function $\widehat T_S+c$ vanishes
identically on no component, and its zero set on each component is discrete.
\item[(c)] For every $H\in\DA(S)$, every $v\in S$ and every $c\in\R$,
\begin{equation}
\label{eq:plancherel}
  W_v(H)=\int_{\widehat{\CS}}\widehat H(\chi)\,\widehat W_v(\chi)\,d\chi,\qquad
  T^{(c)}_S(H)=\int_{\widehat{\CS}}\widehat H(\chi)
     \bigl(\widehat T_S(\chi)+c\bigr)\,d\chi ,
\end{equation}
the integrals being absolutely convergent.
\item[(d)] The function $\widehat T_S(\cz,\cdot)$ is even, and there is a
constant $C$ with $|\widehat T_S(\cz,\tau)|\le C\log(2+|\tau|)$ for all real
$\tau$.
\end{enumerate}
\end{proposition}

Part (a) is contained in \cite[Theorem~5.12(a),(b)]{P43}, part (b) is
\cite[Theorem~5.12(c)]{P43}, and part (c) is \cite[Theorem~5.12(d)]{P43}.
In part (d) the evenness follows from (a), since replacing $\tau$ by
$-\tau$ conjugates $p^{-1/2+i\tau}$ and $\psi(\tfrac14+\tfrac{i\tau}2)$,
and the bound is \cite[(67)]{P43} on $\cz$.

\begin{proposition}[{\cite[Theorem~5.11]{P43}}]
\label{prop:tate}
The functionals
\begin{equation}
\label{eq:tatefunctionals}
  \lambda_\mp(F):=\int_{\CS}F(u)\,|u|_S^{\mp1/2}\,d^*u
\end{equation}
on $\DA(S)$ are multiplicative for $\star$ and satisfy
$\lambda_\mp(F^\dagger)=\overline{\lambda_\pm(F)}$.  The map
$M:=(\lambda_-,\lambda_+)$ is a surjective homomorphism of algebras from
$\DA(S)$ to $\C^2$.  On $\C^2$, with coordinates $z=(z_-,z_+)$, we use the
nondegenerate Hermitian form
\begin{equation}
\label{eq:betaT}
  \beta_T(z,w)=z_-\overline{w_+}+z_+\overline{w_-} .
\end{equation}
\end{proposition}

The notation $\beta_T$ is that of \cite{P43}; the subscript refers to Tate
and not to the trace.  As in \cite[Definition~6.6]{P43}, an element
$F\in\DA(S)$ is called \emph{primitive} if $M(F)=0$.

We use the following conventions on $\R$.  For $u,g\in C_c^\infty(\R)$ we
write $u\star g$ for the additive convolution, $u^\dagger(t)=\overline{u(-t)}$
and $u^-(t)=u(-t)$, and
\begin{equation}
\label{eq:conventions}
  \hat u(\tau)=\int_\R u(t)e^{-i\tau t}\,dt,\qquad
  \Phi_u(s)=\int_\R u(t)e^{(s-\frac12)t}\,dt,\qquad
  \Phi_u(\tfrac12+i\tau)=\hat u(-\tau).
\end{equation}

\begin{proposition}[{\cite[Definitions~3.2 and~3.6, Lemma~3.7]{P44}}]
\label{prop:radalg}
For $h\in C_c^\infty(\R)$ put
\[
  \Gamma_h(w):=h\bigl(-\log|w|_S\bigr),\qquad w\in\CS ,
\]
and let $\DAr(S):=\{\Gamma_h:h\in C_c^\infty(\R)\}$, the set of \emph{radial
classes}.  Then $\DAr(S)$ is the set of elements of $\DA(S)$ that are
invariant under $K_S$, and for $h,g\in C_c^\infty(\R)$
\begin{equation}
\label{eq:radalg}
  \Gamma_h\star\Gamma_g=\Gamma_{h\star g},\qquad
  \Gamma_h^\dagger=\Gamma_{h^\dagger},\qquad
  \lambda_\mp(\Gamma_h)=\int_\R h(t)e^{\pm t/2}\,dt .
\end{equation}
In particular $\DAr(S)$ is a $\dagger$-subalgebra of $\DA(S)$.  The
transform $\widehat{\Gamma_h}$ vanishes off the principal component, and
\begin{equation}
\label{eq:profile}
  \widehat{\Gamma_h}(\cz,\tau)=\hat h(-\tau)=\Phi_h(\tfrac12+i\tau).
\end{equation}
\end{proposition}

In \cite{P44} the group $K_S$ is denoted $K^\natural$ and $h^\dagger$ is
denoted $h^\vee$.  In the coordinate $t$ of \eqref{eq:disint} the function
$\Gamma_h$ is $t\mapsto h(-t)$; in particular $\Gamma_h(1)=h(0)$, and
$h\mapsto\Gamma_h$ is injective.  By
\eqref{eq:radalg} and \eqref{eq:conventions},
\begin{equation}
\label{eq:tatePhi}
  \lambda_-(\Gamma_h)=\Phi_h(1),\qquad\lambda_+(\Gamma_h)=\Phi_h(0),
\end{equation}
so that $\Gamma_h$ is primitive if and only if $\Phi_h(0)=\Phi_h(1)=0$; we
then also call $h$ primitive.  If $F\in C_c^\infty(\R)$ is real, the numbers
\eqref{eq:tatePhi} for $h=F$ are real, and \eqref{eq:betaT} gives
\begin{equation}
\label{eq:polarterm}
  \beta_T\bigl(M(\Gamma_F),M(\Gamma_F)\bigr)=2\,\Phi_F(0)\,\Phi_F(1).
\end{equation}

We turn to the change of chart.  For charts $S\subseteq S'$ let
$\Inf_{SS'}\colon\DA(S)\to\DA(S')$ be the composition with the coordinate
projection $C_{S'}\to\CS$; it is an injective homomorphism of
$\dagger$-algebras \cite[Lemma~8.2]{P43}.  Let $\DAl:=\varinjlim_S\DA(S)$ be
the union of the algebras $\DA(S)$ along these maps.  A radial class is the
same function of $|w|_S$ in every chart, $\Inf_{SS'}\Gamma_h=\Gamma_h$
\cite[Lemma~3.7(v)]{P44}, and so defines an element of $\DAl$, again denoted
$\Gamma_h$.  For $F\in\DA(S)$ let $\rd(F)$ be the supremum of
$\bigl|\log|w|_S\bigr|$ over $w\in\supp F$, with $\rd(0)=0$; thus
\[
  \rd(\Gamma_h)=\sup\{|t|:t\in\supp h\}\qquad(h\ne0).
\]
We write $\ell_S$ for the least prime not in $S_f$.

\begin{proposition}[{\cite[Theorems~8.3 and~8.4]{P43}}]
\label{prop:colimit}
Let $c\in\R$ and $F\in\DA(S)$.  The value $T^{(c)}_{S'}(\Inf_{SS'}F)$ is the
same for all charts $S'$ that contain $S$ and every prime $\ell$ with
$\log\ell\le\rd(F)$.  Denoting it by $T^{(c)}_{\lim}(F)$, one obtains a
linear functional $T^{(c)}_{\lim}$ on $\DAl$, the form
$I^{(c)}_{\lim}(F,G):=T^{(c)}_{\lim}(F\star G^\dagger)$ is Hermitian, and
$T^{(c)}_{\lim}(F)=T^{(c)}_S(F)$ whenever $\rd(F)<\log\ell_S$.  We write
$\Tl:=T^{(0)}_{\lim}$ and $\Il:=I^{(0)}_{\lim}$.  For $h\in C_c^\infty(\R)$,
\begin{equation}
\label{eq:chartchange}
  T^{(c)}_{\lim}(\Gamma_h)=T^{(c)}_S(\Gamma_h)
   +\sum_{\ell\notin S_f}\sum_{k\ge1}(\log\ell)\,\ell^{-k/2}
     \bigl[h(k\log\ell)+h(-k\log\ell)\bigr],
\end{equation}
where $\ell$ runs over the primes not in $S_f$ and only the primes with
$\log\ell\le\rd(\Gamma_h)$ contribute.
\end{proposition}

Formula \eqref{eq:chartchange} is \cite[(98)]{P43} for $F=\Gamma_h$: the term
of the prime $\ell\notin S_f$ there is
$(\log\ell)\sum_{k\ne0}\ell^{-|k|/2}\Gamma_h(w_{\ell,k})$ with
$w_{\ell,k}\in\CS$ of module $\ell^{-k}$, so that
$\Gamma_h(w_{\ell,k})=h(k\log\ell)$.

\begin{lemma}
\label{lem:radial-trace}
Let $h\in C_c^\infty(\R)$.  Then
\begin{equation}
\label{eq:radialtrace}
  T_S(\Gamma_h)=\sum_{p\in S_f}\sum_{k\ge1}(\log p)\,p^{-k/2}
     \bigl[h(k\log p)+h(-k\log p)\bigr]
   +\frac1{2\pi}\int_\R\hat h(-\tau)\,\widehat W_\infty(0,\tau)\,d\tau ,
\end{equation}
where the sum is finite and the integral, which equals
$W_\infty(\Gamma_h)$, converges absolutely.  Moreover, with $\Lambda$ the
von Mangoldt function,
\begin{equation}
\label{eq:radialtracelim}
  \Tl(\Gamma_h)=\sum_{n\ge2}\frac{\Lambda(n)}{\sqrt n}
     \bigl[h(\log n)+h(-\log n)\bigr]
   +\frac1{2\pi}\int_\R\hat h(-\tau)\,\widehat W_\infty(0,\tau)\,d\tau ,
\end{equation}
the sum being finite.  For every $c\in\R$ one has
$T^{(c)}_S(\Gamma_h)=T_S(\Gamma_h)+c\,h(0)$ and
$T^{(c)}_{\lim}(\Gamma_h)=\Tl(\Gamma_h)+c\,h(0)$.
\end{lemma}

\begin{proof}
By \cite[Theorem~3.8]{P44}, for every $p\in S_f$,
\[
  W_p(\Gamma_h)=\sum_{k\ge1}(\log p)\,p^{-k/2}
  \bigl[h(k\log p)+h(-k\log p)\bigr],
\]
a finite sum because $h$ has compact support.  By \eqref{eq:plancherel} and
\eqref{eq:profile},
\[
  W_\infty(\Gamma_h)=\int_{\widehat{\CS}}\widehat{\Gamma_h}\,
  \widehat W_\infty\,d\chi
  =\frac1{2\pi}\int_\R\hat h(-\tau)\,\widehat W_\infty(0,\tau)\,d\tau ,
\]
absolutely convergent.  Summing over $v\in S$ gives \eqref{eq:radialtrace}.
Adding the terms of \eqref{eq:chartchange} extends the sum to all primes, and
since $\Lambda(p^k)=\log p$ and $\Lambda(n)=0$ when $n$ is not a prime power,
the result is \eqref{eq:radialtracelim}.  The last assertion follows from
\eqref{eq:tracec} and $\Gamma_h(1)=h(0)$.
\end{proof}

\section{The explicit formula}
\label{sec:EF}

Throughout, $\xi(s)=\tfrac12s(s-1)\pi^{-s/2}\Gamma(s/2)\zeta(s)$.  It is an
entire function with $\xi(1-s)=\xi(s)$ and $\xi(\bar s)=\overline{\xi(s)}$,
and its zeros, which are the zeros of $\zeta$ in the strip
$0\le\Rel s\le1$, have the same multiplicities as zeros of $\zeta$
\cite[ch.~8 and~12]{Davenport}.  A sum $\sum_\rho$ runs over the zeros of
$\xi$, each repeated according to its multiplicity.  Since $\zeta$ has a
simple pole of residue $1$ at $s=1$ and $\Gamma(\frac12)=\sqrt\pi$, one has
$\xi(1)=\xi(0)=\tfrac12$.

No zero of $\xi$ is real.  Indeed, a real zero would lie in $[0,1]$, and
$\xi(0)=\xi(1)=\tfrac12$.  For $\Rel s>1$ one has
$(1-2^{1-s})\zeta(s)=\sum_{k\ge1}\bigl[(2k-1)^{-s}-(2k)^{-s}\bigr]$, and
since $|(2k-1)^{-s}-(2k)^{-s}|\le|s|\,(2k-1)^{-\Rel s-1}$, the series
converges locally uniformly in $\Rel s>0$ and the identity persists there.
For $0<\sigma<1$ its terms are positive and $1-2^{1-\sigma}<0$, so
$\zeta(\sigma)<0$, and the other factors of $\xi(\sigma)$ do not vanish.

For $b\in\R$ we write $L_b$ for the line $\Rel s=b$ oriented upwards, so
that $\int_{L_b}\vartheta(s)\,ds=i\int_\R\vartheta(b+i\tau)\,d\tau$.

\begin{lemma}
\label{lem:testfn}
Let $h\in C_c^\infty(\R)$ with $\supp h\subseteq[-r,r]$.
\begin{enumerate}
\item[(i)] $\Phi_h$ is entire, and for every integer $N\ge0$ and every
$s=\sigma+i\tau\ne\tfrac12$,
\begin{equation}
\label{eq:Phibound}
  |\Phi_h(s)|\le\|h^{(N)}\|_{L^1}\,e^{r|\sigma-\frac12|}\,|s-\tfrac12|^{-N} .
\end{equation}
Consequently, for every $b_1<b_2$ and every $N\ge0$ there is a constant $C$
such that
\begin{equation}
\label{eq:Phidecay}
  |\Phi_h(s)|\le C\,(1+|s|)^{-N}\qquad(b_1\le\Rel s\le b_2).
\end{equation}
\item[(ii)] $\Phi_h(1-s)=\Phi_{h^-}(s)$, and $\Phi_{h\star g}=\Phi_h\Phi_g$
for $g\in C_c^\infty(\R)$.  In particular $\Phi_h(1-s)=\Phi_h(s)$ if $h$ is
even.
\item[(iii)] Let $F\in C_c^\infty(\R)$ be real with $\supp F\subseteq[-r,r]$,
and put $h:=F\star F^\dagger$.  Then $h$ is real and even,
$\supp h\subseteq[-2r,2r]$, $h(0)=\int_\R F(t)^2\,dt$,
$\hat h=|\hat F|^2$, and
\[
  \Phi_h(s)=\Phi_F(s)\,\Phi_F(1-s),\qquad
  \Phi_h(0)=\Phi_h(1)=\Phi_F(0)\,\Phi_F(1).
\]
\end{enumerate}
\end{lemma}

\begin{proof}
(i) The integral defining $\Phi_h$ extends over a compact set and its
integrand is entire in $s$, so $\Phi_h$ is entire.  Integrating by parts $N$
times, with no boundary terms,
\[
  \Phi_h(s)=(\tfrac12-s)^{-N}\int_\R h^{(N)}(t)\,e^{(s-\frac12)t}\,dt ,
\]
and $|e^{(s-\frac12)t}|=e^{(\sigma-\frac12)t}\le e^{r|\sigma-\frac12|}$ for
$|t|\le r$.  The bound \eqref{eq:Phidecay} follows from \eqref{eq:Phibound}
for the exponents $0$ and $N$.

(ii) The substitution $t\mapsto-t$ gives
$\Phi_h(1-s)=\int h(t)e^{(\frac12-s)t}dt=\int h(-t)e^{(s-\frac12)t}dt$.  By
Fubini's theorem,
\[
  \Phi_{h\star g}(s)=\iint h(u)\,g(t-u)\,e^{(s-\frac12)u}e^{(s-\frac12)(t-u)}
  \,du\,dt=\Phi_h(s)\,\Phi_g(s).
\]

(iii) Since $F$ is real, $F^\dagger=F^-$ and
$h(t)=\int F(u)F(u-t)\,du$, which is real, satisfies $h(-t)=h(t)$ by the
substitution $u\mapsto u+t$, vanishes for $|t|>2r$, and equals
$\int F^2$ at $t=0$.  By (ii),
$\Phi_h(s)=\Phi_F(s)\Phi_{F^-}(s)=\Phi_F(s)\Phi_F(1-s)$.  For real $\tau$,
by \eqref{eq:conventions},
$\hat h(-\tau)=\Phi_h(\tfrac12+i\tau)=\hat F(-\tau)\hat F(\tau)
=|\hat F(\tau)|^2$, because $\hat F(-\tau)=\overline{\hat F(\tau)}$ for real
$F$; and $|\hat F|^2$ is even.
\end{proof}

We use three classical facts about the zeros \cite[ch.~12 and~15]{Davenport}:
the series $\sum_\rho|\rho|^{-2}$ converges; there are constants $C_0$ and
$t_0\ge2$ such that for every $X\ge t_0$ at most $C_0\log X$ zeros $\rho$
satisfy $|\Iml\rho-X|<1$; and, uniformly for $-1\le\sigma\le2$ and
$y\ge t_0$ with $\zeta(\sigma+iy)\ne0$,
\begin{equation}
\label{eq:partialfractions}
  \frac{\zeta'}{\zeta}(\sigma+iy)
  =\sum_{|\Iml\rho-y|<1}\frac1{\sigma+iy-\rho}+O(\log y).
\end{equation}
Logarithmic differentiation of the definition of $\xi$ gives the identity of
meromorphic functions
\begin{equation}
\label{eq:xilog}
  \frac{\xi'}{\xi}(s)=A_0(s)-\tfrac12\log\pi+\tfrac12\psi(s/2)
   +\frac{\zeta'}{\zeta}(s),
  \qquad A_0(s):=\frac1s+\frac1{s-1},
\end{equation}
and $\xi(1-s)=\xi(s)$ gives
\begin{equation}
\label{eq:xireflect}
  \frac{\xi'}{\xi}(1-s)=-\frac{\xi'}{\xi}(s),\qquad A_0(1-s)=-A_0(s).
\end{equation}
For $\Rel s>1$ one has
\begin{equation}
\label{eq:euler}
  -\frac{\zeta'}{\zeta}(s)=\sum_{n\ge2}\Lambda(n)\,n^{-s}
  =\sum_p\frac{(\log p)\,p^{-s}}{1-p^{-s}},
\end{equation}
both series converging absolutely.  Finally, $\psi$ is holomorphic in
$\Rel z>0$ and $\psi(z)=\log z+O(1/|z|)$ as $|z|\to\infty$ in this
half-plane \cite[\S5.11]{NIST}; consequently, for every $\sigma_0>0$,
\begin{equation}
\label{eq:psibound}
  \psi(s/2)=O\bigl(\log(2+|s|)\bigr)\qquad(\Rel s\ge\sigma_0).
\end{equation}

\begin{lemma}
\label{lem:contour}
Let $\varphi$ be an entire function such that
\begin{equation}
\label{eq:gdecay}
  |\varphi(s)|\le C\,(1+|s|)^{-2}\qquad(-1\le\Rel s\le2)
\end{equation}
for some constant $C$.  Then $\sum_\rho \varphi(\rho)$ converges absolutely.  If
moreover $\varphi(1-s)=\varphi(s)$, then the
functions $\varphi \,\xi'/\xi$ and $\varphi A_0$ are integrable on $L_2$, and
\begin{equation}
\label{eq:zeroint}
  \sum_\rho \varphi(\rho)=\frac1{\pi i}\int_{L_2}\varphi(s)\,\frac{\xi'}{\xi}(s)\,ds ,
  \qquad
  \varphi(0)+\varphi(1)=\frac1{\pi i}\int_{L_2}\varphi(s)\,A_0(s)\,ds .
\end{equation}
\end{lemma}

\begin{proof}
Since the zeros lie in $0\le\Rel s\le1$, \eqref{eq:gdecay} gives
$\sum_\rho|\varphi(\rho)|\le C\sum_\rho|\rho|^{-2}<\infty$.  Assume from now on
that $\varphi(1-s)=\varphi(s)$.

Let $\vartheta$ be either $\xi'/\xi$ or $A_0$, so that $\vartheta(1-s)=-\vartheta(s)$ by
\eqref{eq:xireflect}.  On $L_2$ the function $A_0$ is bounded,
$|\zeta'/\zeta(s)|\le\sum_{n\ge2}\Lambda(n)n^{-2}<\infty$, and
$\psi(s/2)=O(\log(2+|s|))$ by \eqref{eq:psibound}; so by \eqref{eq:xilog} and
\eqref{eq:gdecay} the function $\varphi \vartheta$ is $O\bigl(|s|^{-2}\log(2+|s|)\bigr)$ on
$L_2$, hence integrable there.  Since $\varphi(1-s)\vartheta(1-s)=-\varphi(s)\vartheta(s)$, the same
bound holds on $L_{-1}$.  The substitution $s=1-u$ maps $L_{-1}$, oriented
upwards, to $L_2$, oriented downwards, and $ds=-du$; hence
\begin{equation}
\label{eq:reflect}
  \int_{L_{-1}}\varphi(s)\vartheta(s)\,ds=\int_{L_2}\varphi(1-u)\vartheta(1-u)\,du
  =-\int_{L_2}\varphi(u)\vartheta(u)\,du .
\end{equation}

Next we choose heights.  Let $j\ge t_0$ be an integer.  The zeros with
$j\le\Iml\rho\le j+1$ satisfy $|\Iml\rho-(j+\tfrac12)|<1$, so there are at
most $n_j:=\lfloor C_0\log(j+1)\rfloor$ of them.  Divide $[j,j+1]$ into
$n_j+1$ intervals of equal length.  The interior of one of them contains the
ordinate of no zero; let $X_j$ be its midpoint.  Then
\begin{equation}
\label{eq:heights}
  |\Iml\rho-X_j|\ge\frac1{2(n_j+1)}\qquad\text{for every zero }\rho .
\end{equation}
For $-1\le\sigma\le2$ each term of the sum in \eqref{eq:partialfractions} at
$y=X_j$ is at most $2(n_j+1)$ in absolute value, and there are at most
$C_0\log(j+1)$ terms, so $\zeta'/\zeta(\sigma+iX_j)=O(\log^2j)$.  For
$\tfrac12\le\sigma\le2$ the functions $A_0$ and $\psi(s/2)$ are $O(\log j)$
at $s=\sigma+iX_j$, by \eqref{eq:psibound}, so \eqref{eq:xilog} gives
$\xi'/\xi(\sigma+iX_j)=O(\log^2j)$.  By \eqref{eq:xireflect} and
$\xi'/\xi(\bar s)=\overline{\xi'/\xi(s)}$,
\[
  \Bigl|\frac{\xi'}{\xi}(\sigma+iX_j)\Bigr|
  =\Bigl|\frac{\xi'}{\xi}(1-\sigma+iX_j)\Bigr| ,
\]
so the bound extends to $-1\le\sigma\le\tfrac12$, and by conjugation to the
height $-X_j$:
\begin{equation}
\label{eq:horizontal}
  \frac{\xi'}{\xi}(\sigma\pm iX_j)=O(\log^2j)
  \qquad\text{uniformly for }-1\le\sigma\le2 .
\end{equation}
The function $A_0$ is bounded by $2$ on these segments.

Let $Q_j$ be the rectangle with vertices $-1\mp iX_j$, $2\mp iX_j$.  The
zeros of $\xi$ are symmetric under $\rho\mapsto\bar\rho$, so by
\eqref{eq:heights} none has ordinate $\pm X_j$; and none lies on $L_2$ or
$L_{-1}$.  Thus $\varphi \vartheta$ is holomorphic on a neighbourhood of $\partial Q_j$.
The function $\xi'/\xi$ has a simple pole at each zero of $\xi$, with residue
the multiplicity, and no other pole, while $A_0$ has simple poles of residue
$1$ at $s=0$ and $s=1$; and $\varphi$ is entire.  By the residue theorem, with
$\partial Q_j$ traversed counterclockwise,
\[
  \frac1{2\pi i}\oint_{\partial Q_j}\varphi \,\frac{\xi'}{\xi}\,ds
  =\sum_{|\Iml\rho|<X_j}\varphi(\rho),\qquad
  \frac1{2\pi i}\oint_{\partial Q_j}\varphi \,A_0\,ds=\varphi(0)+\varphi(1).
\]
The two horizontal sides have length $3$, and on them $\varphi \vartheta$ is
$O(X_j^{-2}\log^2j)$ by \eqref{eq:gdecay} and \eqref{eq:horizontal}; their
contribution tends to $0$ as $j\to\infty$.  The integrals over the vertical
sides tend to the integrals over $L_2$ and $L_{-1}$.  Hence
\[
  \frac1{2\pi i}\Bigl[\int_{L_2}-\int_{L_{-1}}\Bigr]\varphi(s)\vartheta(s)\,ds
\]
is equal to $\sum_\rho \varphi(\rho)$ for $\vartheta=\xi'/\xi$ and to $\varphi(0)+\varphi(1)$ for
$\vartheta=A_0$.  By \eqref{eq:reflect} the integral over $L_{-1}$ is the negative of
the integral over $L_2$.  This is \eqref{eq:zeroint}.
\end{proof}

\begin{theorem}
\label{thm:EF}
Let $h\in C_c^\infty(\R)$ with $\supp h\subseteq[-r,r]$.  Then
$\sum_\rho\Phi_h(\rho)$ converges absolutely, and
\begin{equation}
\label{eq:EFh}
\begin{aligned}
  \sum_\rho\Phi_h(\rho)={}&\Phi_h(0)+\Phi_h(1)
   +\frac1{2\pi}\int_\R\hat h(\tau)
     \bigl[\Rel\psi(\tfrac14+\tfrac{i\tau}2)-\log\pi\bigr]\,d\tau\\
   &-\sum_{2\le n<e^{r}}\frac{\Lambda(n)}{\sqrt n}
      \bigl[h(\log n)+h(-\log n)\bigr],
\end{aligned}
\end{equation}
the integral being absolutely convergent.  In particular, if
$F\in C_c^\infty(\R)$ is real with $\supp F\subseteq[-r,r]$, then
\begin{equation}
\label{eq:EF}
\begin{aligned}
  \sum_\rho\Phi_F(\rho)\Phi_F(1-\rho)
  ={}&2\,\Phi_F(0)\Phi_F(1)
  +\frac1{2\pi}\int_\R|\hat F(\tau)|^2
     \bigl[\Rel\psi(\tfrac14+\tfrac{i\tau}2)-\log\pi\bigr]\,d\tau\\
  &-\sum_{2\le n<e^{2r}}\frac{\Lambda(n)}{\sqrt n}
     \bigl[h(\log n)+h(-\log n)\bigr],
\end{aligned}
\end{equation}
with $h=F\star F^\dagger$.
\end{theorem}

\begin{proof}
The function $\hat h$ is a Schwartz function and
$\Rel\psi(\tfrac14+\tfrac{i\tau}2)=O(\log(2+|\tau|))$ by
\eqref{eq:psibound}, so the integral converges absolutely.  By
\eqref{eq:Phidecay} the function $\Phi_h$ satisfies \eqref{eq:gdecay}; hence
$\sum_\rho|\Phi_h(\rho)|<\infty$ by the first assertion of
Lemma~\ref{lem:contour}.

We first reduce to even $h$.  Let $h_e:=\tfrac12(h+h^-)$.  By
Lemma~\ref{lem:testfn}(ii),
$\Phi_{h_e}(s)=\tfrac12\bigl(\Phi_h(s)+\Phi_h(1-s)\bigr)$.  Since
$\xi(1-s)=\xi(s)$, the map $\rho\mapsto1-\rho$ is a bijection of the zeros
preserving multiplicities, so
$\sum_\rho\Phi_h(1-\rho)=\sum_\rho\Phi_h(\rho)$ and the left-hand side of
\eqref{eq:EFh} is the same for $h$ and for $h_e$.  So is the right-hand
side: $\Phi_{h_e}(0)+\Phi_{h_e}(1)=\Phi_h(0)+\Phi_h(1)$; the function
$\tau\mapsto\Rel\psi(\tfrac14+\tfrac{i\tau}2)$ is even, because
$\psi(\bar z)=\overline{\psi(z)}$, and
$\widehat{h_e}(\tau)=\tfrac12(\hat h(\tau)+\hat h(-\tau))$; and
$h_e(t)+h_e(-t)=h(t)+h(-t)$.  We may therefore assume that $h$ is even.

Let $h$ be even and put $\varphi:=\Phi_h$.  Then $\varphi$ is entire, $\varphi(1-s)=\varphi(s)$ by
Lemma~\ref{lem:testfn}(ii), and by \eqref{eq:conventions}
\begin{equation}
\label{eq:gline}
  \varphi(\tfrac12+i\tau)=\hat h(-\tau)=\hat h(\tau)\qquad(\tau\in\R).
\end{equation}
By Lemma~\ref{lem:contour} and \eqref{eq:xilog},
\begin{equation}
\label{eq:threeterms}
  \sum_\rho \varphi(\rho)=\varphi(0)+\varphi(1)
  +\frac1{\pi i}\int_{L_2}\varphi(s)\bigl[-\tfrac12\log\pi+\tfrac12\psi(s/2)\bigr]ds
  +\frac1{\pi i}\int_{L_2}\varphi \,\frac{\zeta'}{\zeta}\,ds ,
\end{equation}
each integrand being integrable on $L_2$, as shown in the proof of
Lemma~\ref{lem:contour}.  We evaluate the two integrals.

The function $s\mapsto-\tfrac12\log\pi+\tfrac12\psi(s/2)$ is holomorphic in
$\Rel s>0$ and is $O(\log(2+|s|))$ for $\Rel s\ge\tfrac12$ by
\eqref{eq:psibound}.  For $X>0$, Cauchy's theorem on the rectangle with
vertices $\tfrac12\mp iX$, $2\mp iX$, whose horizontal sides contribute
$O(X^{-2}\log(2+X))$ by \eqref{eq:Phidecay}, gives, on letting
$X\to\infty$, that the integral over $L_2$ equals the integral over
$L_{1/2}$; both integrals converge absolutely by \eqref{eq:Phidecay} and
\eqref{eq:psibound}.  With $s=\tfrac12+i\tau$ and
\eqref{eq:gline},
\[
  \frac1{\pi i}\int_{L_2}\varphi(s)\bigl[-\tfrac12\log\pi+\tfrac12\psi(s/2)\bigr]ds
  =\frac1{2\pi}\int_\R\hat h(\tau)
   \bigl[\psi(\tfrac14+\tfrac{i\tau}2)-\log\pi\bigr]\,d\tau .
\]
The function $\tau\mapsto\Iml\psi(\tfrac14+\tfrac{i\tau}2)$ is odd, while
$\hat h$ is even; so $\psi$ may be replaced by $\Rel\psi$, and the right-hand
side is the integral in \eqref{eq:EFh}.

On $L_2$ one has $\zeta'/\zeta(s)=-\sum_{n\ge2}\Lambda(n)n^{-s}$ with
$\sum_n\Lambda(n)n^{-2}<\infty$, and $\varphi$ is integrable on $L_2$; so
\[
  \frac1{\pi i}\int_{L_2}\varphi \,\frac{\zeta'}{\zeta}\,ds
  =-\sum_{n\ge2}\Lambda(n)\,\frac1{\pi i}\int_{L_2}\varphi(s)\,n^{-s}\,ds .
\]
For each $n$ the function $\varphi(s)n^{-s}$ is entire, and $|n^{-s}|\le1$ for
$\Rel s\ge\tfrac12$.  By the same argument on the same rectangles, whose horizontal sides now
contribute $O(X^{-2})$, the line $L_2$ may be replaced by $L_{1/2}$, and by \eqref{eq:gline} and Fourier inversion
\[
  \frac1{\pi i}\int_{L_2}\varphi(s)\,n^{-s}\,ds
  =\frac{2}{\sqrt n}\cdot\frac1{2\pi}\int_\R\hat h(\tau)\,
   e^{-i\tau\log n}\,d\tau
  =\frac{2}{\sqrt n}\,h(-\log n)
  =\frac1{\sqrt n}\bigl[h(\log n)+h(-\log n)\bigr],
\]
because $h$ is even.  This vanishes for $\log n\ge r$, since $h$ is
continuous with support in $[-r,r]$.

Inserting the two values into \eqref{eq:threeterms} gives \eqref{eq:EFh}.
Formula \eqref{eq:EF} is \eqref{eq:EFh} for $h=F\star F^\dagger$, whose
support lies in $[-2r,2r]$, by Lemma~\ref{lem:testfn}(iii).
\end{proof}

Formula \eqref{eq:EFh} is the explicit formula of Weil \cite{Weil1952} for
the test function $h$; see also \cite{Barner1981}.

\section{The comparison theorem}
\label{sec:comparison}

\begin{theorem}
\label{thm:comparison}
\begin{enumerate}
\item[(i)] For every $h\in C_c^\infty(\R)$,
\begin{equation}
\label{eq:comparisonh}
  \Tl(\Gamma_h)=\Phi_h(0)+\Phi_h(1)-\sum_\rho\Phi_h(\rho),
\end{equation}
the sum being absolutely convergent.
\item[(ii)] For every real $F\in C_c^\infty(\R)$,
\begin{equation}
\label{eq:comparison}
\begin{aligned}
  \Il(\Gamma_F,\Gamma_F)=\Tl(\Gamma_{F\star F^\dagger})
  &=2\,\Phi_F(0)\Phi_F(1)-\sum_\rho\Phi_F(\rho)\Phi_F(1-\rho)\\
  &=\beta_T\bigl(M(\Gamma_F),M(\Gamma_F)\bigr)
     -\sum_\rho\Phi_F(\rho)\Phi_F(1-\rho).
\end{aligned}
\end{equation}
In particular, if $F$ is primitive, then
\begin{equation}
\label{eq:comparisonprim}
  \Il(\Gamma_F,\Gamma_F)=-\sum_\rho\Phi_F(\rho)\Phi_F(1-\rho).
\end{equation}
\item[(iii)] For every chart $S$ and every $h\in C_c^\infty(\R)$,
\begin{equation}
\label{eq:comparisonS}
  T_S(\Gamma_h)=\Phi_h(0)+\Phi_h(1)-\sum_\rho\Phi_h(\rho)
   -\sum_{\ell\notin S_f}\sum_{k\ge1}(\log\ell)\,\ell^{-k/2}
     \bigl[h(k\log\ell)+h(-k\log\ell)\bigr],
\end{equation}
where $\ell$ runs over the primes not in $S_f$ and the double sum is finite.
The double sum vanishes if $h$ vanishes outside the interval
$(-\log\ell_S,\log\ell_S)$.  In that case $T_S(\Gamma_h)$ is equal to the
right-hand side of \eqref{eq:comparisonh}; and for real $F$ the number
$I_S(\Gamma_F,\Gamma_F)$ is equal to the right-hand side of
\eqref{eq:comparison} if $F\star F^\dagger$ vanishes outside that interval.
\end{enumerate}
\end{theorem}

\begin{proof}
(i) By \eqref{eq:radialtracelim} and \eqref{eq:symbolclosed},
\[
  \Tl(\Gamma_h)=\sum_{n\ge2}\frac{\Lambda(n)}{\sqrt n}
     \bigl[h(\log n)+h(-\log n)\bigr]
   +\frac1{2\pi}\int_\R\hat h(-\tau)
     \bigl[\log\pi-\Rel\psi(\tfrac14+\tfrac{i\tau}2)\bigr]\,d\tau .
\]
The bracket is even in $\tau$, so $\hat h(-\tau)$ may be replaced by
$\hat h(\tau)$.  The sum extends over $n<e^{r}$ if $\supp h\subseteq[-r,r]$.
Comparing with \eqref{eq:EFh} gives \eqref{eq:comparisonh}.

(ii) By \eqref{eq:radalg},
$\Gamma_F\star\Gamma_F^\dagger=\Gamma_{F\star F^\dagger}$, so
$\Il(\Gamma_F,\Gamma_F)=\Tl(\Gamma_h)$ with $h=F\star F^\dagger$.  By
Lemma~\ref{lem:testfn}(iii), $\Phi_h(0)+\Phi_h(1)=2\Phi_F(0)\Phi_F(1)$ and
$\Phi_h(\rho)=\Phi_F(\rho)\Phi_F(1-\rho)$, so the first line of
\eqref{eq:comparison} is \eqref{eq:comparisonh}.  The second line is
\eqref{eq:polarterm}, and it vanishes for primitive $F$.

(iii) Formula \eqref{eq:comparisonS} is \eqref{eq:chartchange} for $c=0$
together with (i).  Every number $\pm k\log\ell$ occurring in the double sum
has absolute value at least $\log\ell_S$.  The last assertion follows from
$I_S(\Gamma_F,\Gamma_F)=T_S(\Gamma_{F\star F^\dagger})$ and the computation
in (ii).
\end{proof}

\begin{corollary}
\label{cor:c-unique}
Let $c\in\R$ and let $F\in C_c^\infty(\R)$ be real.  Then
\begin{equation}
\label{eq:comparisonc}
  I^{(c)}_{\lim}(\Gamma_F,\Gamma_F)
   =2\,\Phi_F(0)\Phi_F(1)-\sum_\rho\Phi_F(\rho)\Phi_F(1-\rho)
   +c\int_\R F(t)^2\,dt .
\end{equation}
Consequently, if $F\ne0$, the identity \eqref{eq:comparison} holds with
$I^{(c)}_{\lim}$ in place of $\Il$ if and only if $c=0$.  Likewise, for
$h\in C_c^\infty(\R)$,
\begin{equation}
\label{eq:comparisonhc}
  T^{(c)}_{\lim}(\Gamma_h)=\Phi_h(0)+\Phi_h(1)-\sum_\rho\Phi_h(\rho)+c\,h(0);
\end{equation}
so, for a given $h$ with $h(0)\ne0$, the identity \eqref{eq:comparisonh}
holds with $T^{(c)}_{\lim}$ in place of $\Tl$ if and only if $c=0$.
\end{corollary}

\begin{proof}
Let $h=F\star F^\dagger$.  By Lemma~\ref{lem:radial-trace},
$I^{(c)}_{\lim}(\Gamma_F,\Gamma_F)=T^{(c)}_{\lim}(\Gamma_h)
=\Tl(\Gamma_h)+c\,h(0)$, and $h(0)=\int F^2$ by
Lemma~\ref{lem:testfn}(iii).  Now apply Theorem~\ref{thm:comparison}(ii).
The assertion on $c$ holds because $\int F^2>0$ for $F\ne0$.  Formula
\eqref{eq:comparisonhc} is Lemma~\ref{lem:radial-trace} and
Theorem~\ref{thm:comparison}(i).
\end{proof}

\section{Residues}
\label{sec:polar}

For a prime $p$ and a chart $S$ we define meromorphic functions on $\C$ by
\begin{equation}
\label{eq:AS}
  f_p(s):=\frac{(\log p)\,p^{-s}}{1-p^{-s}},\qquad
  a(s):=\tfrac12\log\pi-\tfrac12\psi(s/2),\qquad
  A_S(s):=a(s)+\sum_{p\in S_f}f_p(s),
\end{equation}
and
\begin{equation}
\label{eq:A}
  A(s):=a(s)-\frac{\zeta'}{\zeta}(s).
\end{equation}
By \eqref{eq:euler},
\begin{equation}
\label{eq:Asum}
  A(s)=a(s)+\sum_pf_p(s)\qquad(\Rel s>1),
\end{equation}
the sum over all primes converging absolutely.

\begin{lemma}
\label{lem:localsymbols}
Let $S$ be a chart.
\begin{enumerate}
\item[(i)] For real $\tau$ and $s=\tfrac12+i\tau$,
\[
  \widehat W_p(\cz,\tau)=f_p(s)+f_p(1-s)\quad(p\in S_f),\qquad
  \widehat W_\infty(0,\tau)=a(s)+a(1-s),
\]
so that $\widehat T_S(\cz,\tau)=A_S(s)+A_S(1-s)$.
\item[(ii)] The poles of $f_p$ are the points of $(2\pi i/\log p)\Z$ and the
poles of $a$ are the points $-2n$, $n\ge0$; all of them are simple, with
residue $1$.  Hence $A_S$ is holomorphic in $\Rel s>0$, and
$\Res_{s=0}A_S=|S_f|+1$.
\item[(iii)] For $\Rel s=\sigma>0$ one has $|f_p(s)|\le(\log p)/(p^\sigma-1)$,
and for every $\sigma_0>0$ there is a constant $C$ with
$|a(s)|\le C\log(2+|s|)$ for $\Rel s\ge\sigma_0$.  For $\Rel s=\sigma>1$,
\begin{equation}
\label{eq:AminusAS}
  |A(s)-A_S(s)|\le\sum_{p\notin S_f}\frac{\log p}{p^\sigma-1}.
\end{equation}
\end{enumerate}
\end{lemma}

\begin{proof}
(i) Here $1-s=\tfrac12-i\tau=\bar s$.  By \eqref{eq:symbolclosed},
$\widehat W_p(\cz,\tau)=2\Rel f_p(1-s)$, and
$\overline{f_p(1-s)}=f_p(\overline{1-s})=f_p(s)$.  In the same way
$\widehat W_\infty(0,\tau)=\log\pi-\Rel\psi(s/2)$ and
$\overline{\psi(s/2)}=\psi((1-s)/2)$.

(ii) The denominator $1-p^{-s}$ vanishes exactly on $(2\pi i/\log p)\Z$, and
at such a point $s_0$ one has $1-p^{-s}=(s-s_0)\log p+O((s-s_0)^2)$ and
$p^{-s_0}=1$, so the residue of $f_p$ is $1$.  By the expansion
$\psi(z)=-\gamma+\sum_{n\ge0}\bigl[\frac1{n+1}-\frac1{n+z}\bigr]$
\cite[(5.7.6)]{NIST}, the poles of $\psi$ are the points $-n$, $n\ge0$, and
they are simple with residue $-1$; so $-\tfrac12\psi(s/2)$ has simple poles
of residue $1$ at the points $-2n$.

(iii) $|p^{-s}|=p^{-\sigma}$ and $|1-p^{-s}|\ge1-p^{-\sigma}$ give the first
bound, the second is \eqref{eq:psibound}, and \eqref{eq:AminusAS} follows
from \eqref{eq:Asum} and the first bound.
\end{proof}

By (i) and (ii), the function $\tau\mapsto A_S(s)+A_S(1-s)$,
$s=\tfrac12+i\tau$, is the meromorphic continuation of
$\widehat T_S(\cz,\cdot)$ to $\tau\in\C$.  Its poles at $\tau=\pm i/2$, that
is at $s=0$ and $s=1$, have residues $\mp i(|S_f|+1)$ in the variable $\tau$;
this is \cite[Theorem~5.13(c)]{P43}.

\begin{theorem}
\label{thm:residues}
\begin{enumerate}
\item[(i)] One has
\begin{equation}
\label{eq:Axi}
  A(s)=\frac1s+\frac1{s-1}-\frac{\xi'}{\xi}(s),\qquad A(1-s)=-A(s).
\end{equation}
The poles of $A$ are $s=0$ and $s=1$, which are simple with residue $1$, and
the zeros $\rho$ of $\xi$, which are simple poles of $A$ with residue
$-m_\rho$, where $m_\rho$ is the multiplicity of $\rho$.
\item[(ii)] Let $h\in C_c^\infty(\R)$ be even.  For every chart $S$ and every
$b>0$,
\begin{equation}
\label{eq:lineS}
  T_S(\Gamma_h)=\frac1{\pi i}\int_{L_b}\Phi_h(s)\,A_S(s)\,ds ,
\end{equation}
and for every $b>1$,
\begin{equation}
\label{eq:lineinf}
  \Tl(\Gamma_h)=\frac1{\pi i}\int_{L_b}\Phi_h(s)\,A(s)\,ds ,
\end{equation}
all integrals being absolutely convergent.
\item[(iii)] For $h$ as in \emph{(ii)}, $\Tl(\Gamma_h)$ is the sum of the
residues of $\Phi_hA$:
\begin{equation}
\label{eq:residuesum}
  \Tl(\Gamma_h)
  =\Res_{s=0}(\Phi_hA)+\Res_{s=1}(\Phi_hA)+\sum_\rho\Res_{s=\rho}(\Phi_hA)
  =\lambda_+(\Gamma_h)+\lambda_-(\Gamma_h)-\sum_\rho\Phi_h(\rho),
\end{equation}
the sum running over the distinct zeros in the middle and over the zeros
with multiplicity on the right, and both converging absolutely.  The residue
at $s=0$ is $\Phi_h(0)=\lambda_+(\Gamma_h)$, and the residue at $s=1$ is
$\Phi_h(1)=\lambda_-(\Gamma_h)$.  If $h=F\star F^\dagger$ with
$F\in C_c^\infty(\R)$ real, their sum is
$2\Phi_F(0)\Phi_F(1)=\beta_T\bigl(M(\Gamma_F),M(\Gamma_F)\bigr)$.
\end{enumerate}
\end{theorem}

\begin{proof}
(i) The first identity of \eqref{eq:Axi} is \eqref{eq:xilog}, and the second
follows from it by \eqref{eq:xireflect}.  Since $\xi$ is entire and
$\xi(0)=\xi(1)=\tfrac12$, the function $\xi'/\xi$ is holomorphic at $0$ and
at $1$, and its poles are the zeros of $\xi$, each a simple pole with residue
$m_\rho$.

(ii) By Lemma~\ref{lem:testfn}(ii), $\Phi_h(1-s)=\Phi_h(s)$.  By
\eqref{eq:plancherel} with $c=0$, \eqref{eq:profile} and
Lemma~\ref{lem:localsymbols}(i),
\[
  T_S(\Gamma_h)=\frac1{2\pi}\int_\R\Phi_h(\tfrac12+i\tau)\,
     \widehat T_S(\cz,\tau)\,d\tau
  =\frac1{2\pi i}\int_{L_{1/2}}\Phi_h(s)\bigl[A_S(s)+A_S(1-s)\bigr]ds .
\]
By Lemma~\ref{lem:localsymbols}(iii) and \eqref{eq:Phidecay}, each of the two
terms is absolutely integrable on $L_{1/2}$.  The substitution
$s\mapsto1-s$ maps $L_{1/2}$ to itself with the orientation reversed, so
\[
  \int_{L_{1/2}}\Phi_h(s)A_S(1-s)\,ds=\int_{L_{1/2}}\Phi_h(1-s)A_S(s)\,ds
  =\int_{L_{1/2}}\Phi_h(s)A_S(s)\,ds ,
\]
which proves \eqref{eq:lineS} for $b=\tfrac12$.  Let $b>0$.  In the closed
strip between $L_{1/2}$ and $L_b$ the function $\Phi_hA_S$ is holomorphic by
Lemma~\ref{lem:localsymbols}(ii), and it is $O(|\tau|^{-2}\log|\tau|)$ as
$|\tau|\to\infty$, uniformly in $\sigma$, by
Lemma~\ref{lem:localsymbols}(iii) and \eqref{eq:Phidecay}.  Cauchy's theorem
on the rectangles with vertices $\tfrac12\pm iX$, $b\pm iX$, $X\to\infty$,
gives \eqref{eq:lineS}.

Now let $b>1$.  On $L_b$ the function $\zeta'/\zeta$ is bounded by
\eqref{eq:euler}, so $\Phi_hA$ is absolutely integrable there, by
Lemma~\ref{lem:localsymbols}(iii) and \eqref{eq:Phidecay}.  By
Proposition~\ref{prop:colimit}, $\Tl(\Gamma_h)=T_S(\Gamma_h)$ for every chart
$S$ such that $S_f$ contains the primes $\ell$ with
$\log\ell\le\rd(\Gamma_h)$.  For such $S$, \eqref{eq:lineS} and
\eqref{eq:AminusAS} give
\[
  \Bigl|\Tl(\Gamma_h)-\frac1{\pi i}\int_{L_b}\Phi_hA\,ds\Bigr|
  \le\frac1\pi\int_\R|\Phi_h(b+i\tau)|\,d\tau\cdot
  \sum_{p\notin S_f}\frac{\log p}{p^b-1},
\]
and the right-hand side tends to $0$ as $S_f$ increases to the set of all
primes.  The left-hand side does not depend on $S$, so it vanishes.  This
proves \eqref{eq:lineinf}.

(iii) By \eqref{eq:Phidecay} the function $\Phi_h$ satisfies the
hypotheses of Lemma~\ref{lem:contour}.  By \eqref{eq:lineinf} for $b=2$,
\eqref{eq:Axi} and \eqref{eq:zeroint},
\[
  \Tl(\Gamma_h)=\frac1{\pi i}\int_{L_2}\Phi_h\,A_0\,ds
   -\frac1{\pi i}\int_{L_2}\Phi_h\,\frac{\xi'}{\xi}\,ds
  =\Phi_h(0)+\Phi_h(1)-\sum_\rho\Phi_h(\rho),
\]
the series being absolutely convergent.  By (i), the residues of $\Phi_hA$ at
$0$, at $1$ and at a zero $\rho$ are $\Phi_h(0)$, $\Phi_h(1)$ and
$-m_\rho\Phi_h(\rho)$.  The remaining assertions are \eqref{eq:tatePhi},
Lemma~\ref{lem:testfn}(iii) and \eqref{eq:polarterm}.
\end{proof}

The proof of \eqref{eq:residuesum} does not use Theorem~\ref{thm:EF}; for
even $h$ it is a second proof of \eqref{eq:comparisonh}.  For arbitrary
$h\in C_c^\infty(\R)$ let $h_e:=\tfrac12(h+h^-)$ be its even part.  The
right-hand sides of \eqref{eq:radialtrace} and \eqref{eq:radialtracelim}
are the same for $h$ and for $h^-$, because $\widehat W_\infty(0,\cdot)$ is
even; so $T_S(\Gamma_h)=T_S(\Gamma_{h_e})$ and
$\Tl(\Gamma_h)=\Tl(\Gamma_{h_e})$, and the formulas \eqref{eq:lineS},
\eqref{eq:lineinf} hold for $h$ with $\Phi_h(s)$ replaced by
$\Phi_{h_e}(s)=\tfrac12\bigl(\Phi_h(s)+\Phi_h(1-s)\bigr)$.

\begin{proposition}
\label{cor:noncommute}
Let $S_f$ increase to the set of all primes.  Then $A_S\to A$ uniformly on
every half-plane $\Rel s\ge b$ with $b>1$.  On the other hand
$\Res_{s=0}A_S=|S_f|+1\to\infty$, whereas $\Res_{s=0}A=1$; and
$A_S(\sigma)\to+\infty$ for every $\sigma\in(0,1]$, whereas $A$ is
holomorphic on the interval $(0,1)$, with $A(\tfrac12)=0$, and its pole at
$s=1$ is a pole of no $A_S$.  In particular
$\widehat T_S(\cz,0)=2A_S(\tfrac12)\to+\infty$.
\end{proposition}

\begin{proof}
For $\Rel s\ge b$ the right-hand side of \eqref{eq:AminusAS} is at most
$\sum_{p\notin S_f}(\log p)/(p^b-1)$, the tail of a convergent series.  The
residues are given by Lemma~\ref{lem:localsymbols}(ii) and
Theorem~\ref{thm:residues}(i).  By Theorem~\ref{thm:residues}(i) the poles
of $A$ in $\Rel s>0$ are $s=1$ and the zeros of $\xi$, none of which is real
(Section~\ref{sec:EF}); so $A$ is holomorphic on $(0,1)$, and
$A(\tfrac12)=-A(\tfrac12)$ by \eqref{eq:Axi}.  For $0<\sigma\le1<\sigma'$ and every prime
$p$ one has $f_p(\sigma)\ge f_p(1)\ge f_p(\sigma')$, and by \eqref{eq:euler}
\[
  \sum_pf_p(1)\ \ge\ \sum_pf_p(\sigma')=-\frac{\zeta'}{\zeta}(\sigma'),
\]
which tends to $+\infty$ as $\sigma'\downarrow1$, because $\zeta$ has a pole
at $s=1$.  Hence $\sum_pf_p(1)=\infty$ and
$A_S(\sigma)\ge a(\sigma)+\sum_{p\in S_f}f_p(1)\to+\infty$.  The last
assertion follows from Lemma~\ref{lem:localsymbols}(i) at $\tau=0$.
\end{proof}

Thus the passage to the limit over the charts commutes with the integration
over $L_b$ for $b>1$, by \eqref{eq:lineinf}, but not with the continuation to
$\Rel s\le1$: the residue of $A$ at $s=0$ is that of the archimedean term
$a$ alone, since $A-a=-\zeta'/\zeta$ is holomorphic at $s=0$, and the residue
of $A$ at $s=1$ is that of $-\zeta'/\zeta$.

\section{Zeros off the line; nondegeneracy}
\label{sec:sufficiency}

For $F\in C_c^\infty(\R)$ write
\begin{equation}
\label{eq:tilde}
  \widetilde F(z):=\Phi_F(\tfrac12+iz)=\int_\R F(t)e^{izt}\,dt
  \qquad(z\in\C),
\end{equation}
an entire function, and let $\partial:=i\,d/dt$.

\begin{lemma}
\label{lem:Pi}
Let $G\in C_c^\infty(\R)$ and let $q$ be a polynomial.  Then
$\widetilde{q(\partial)G}(z)=q(z)\,\widetilde G(z)$.  In particular, if $q$
is divisible by $\tfrac14+z^2$, then $q(\partial)G$ is primitive.
\end{lemma}

\begin{proof}
One integration by parts gives
$\int iG'(t)e^{izt}dt=z\int G(t)e^{izt}dt$, that is,
$\widetilde{\partial G}=z\widetilde G$; the general case follows by iteration
and linearity.  The points $s=0$ and $s=1$ correspond to $z=\pm i/2$, the
zeros of $\tfrac14+z^2$.
\end{proof}

Let $\mathcal Z$ be the set of $z\in\C$ with $\xi(\tfrac12+iz)=0$, and let
$m(z)$ be the multiplicity of the zero $\rho=\tfrac12+iz$.  Then
$|\Iml z|\le\tfrac12$ on $\mathcal Z$, the set $\mathcal Z$ is invariant
under $z\mapsto-z$ and under $z\mapsto\bar z$, which correspond to
$\rho\mapsto1-\rho$ and to $\rho\mapsto1-\bar\rho$, and $m$ is invariant
under both maps.  Since $1-\rho=\tfrac12-iz$, for every $F\in C_c^\infty(\R)$
\begin{equation}
\label{eq:zsum}
  \sum_\rho\Phi_F(\rho)\Phi_F(1-\rho)
  =\sum_{z\in\mathcal Z}m(z)\,\widetilde F(z)\,\widetilde F(-z),
\end{equation}
and the series converges absolutely, by Lemma~\ref{lem:contour} applied to
the function $s\mapsto\Phi_F(s)\Phi_F(1-s)$, which satisfies
\eqref{eq:gdecay} by \eqref{eq:Phidecay}.  If $F$ is real then
\begin{equation}
\label{eq:realF}
  \widetilde F(-z)=\overline{\widetilde F(\bar z)}\qquad(z\in\C).
\end{equation}

\begin{theorem}
\label{thm:sufficiency}
Suppose that $\xi$ has a zero $\rho_0$ with $\Rel\rho_0\ne\tfrac12$.  Then
there are nonzero real primitive functions $F_+,F_-\in C_c^\infty(\R)$, with
$F_+$ even and $F_-$ odd, such that
\begin{equation}
\label{eq:negative}
  \sum_\rho\Phi_{F_\pm}(\rho)\Phi_{F_\pm}(1-\rho)<0 ,
  \qquad\text{that is,}\qquad \Il(\Gamma_{F_\pm},\Gamma_{F_\pm})>0 .
\end{equation}
Conversely, if every zero of $\xi$ lies on the line $\Rel s=\tfrac12$, then
$\sum_\rho\Phi_F(\rho)\Phi_F(1-\rho)\ge0$ for every real
$F\in C_c^\infty(\R)$.
\end{theorem}

\begin{proof}
We prove the converse first.  If every zero lies on the line then
$\mathcal Z\subset\R$, and for real $F$ and real $z$ the term
$\widetilde F(z)\widetilde F(-z)$ of \eqref{eq:zsum} equals
$|\widetilde F(z)|^2\ge0$ by \eqref{eq:realF}.

Now let $\rho_0$ be a zero with $\Rel\rho_0\ne\tfrac12$, and write
$\rho_0=\tfrac12+iz_0$, $z_0=x_0+iy_0$, so that $x_0=\Iml\rho_0$ and
$y_0=\tfrac12-\Rel\rho_0\ne0$.  No zero of $\xi$ is real, so $x_0\ne0$.
Hence $\Iml(z_0^2)=2x_0y_0\ne0$, the four points
$z_0,-z_0,\bar z_0,-\bar z_0$ are distinct, they all belong to $\mathcal Z$,
and they have the same multiplicity $m_0:=m(z_0)\ge1$.  Let $\mathcal O_0$
be the set of these four points.

\emph{The auxiliary function.}  Fix a real even $\eta_1\in C_c^\infty(\R)$
with $\int\eta_1=1$ and put $\eta_\varepsilon(t):=\varepsilon^{-1}
\eta_1(t/\varepsilon)$.  Then $\widetilde{\eta_\varepsilon}(z_0)
=\widetilde{\eta_1}(\varepsilon z_0)\to\widetilde{\eta_1}(0)=1$ as
$\varepsilon\to0$.  Fix $\varepsilon$ with
$\widetilde{\eta_\varepsilon}(z_0)\ne0$ and put $\eta:=\eta_\varepsilon$,
$u:=\widetilde\eta$.  The function $u$ is entire and even, and
$u(\bar z)=\overline{u(z)}$ by \eqref{eq:realF}.  For $s=\tfrac12+iz$ one
has $1+|s|\ge\tfrac12(1+|z|)$, and $0\le\Rel s\le1$ if
$|\Iml z|\le\tfrac12$.  So by \eqref{eq:Phidecay}, for
every $N\ge0$ there is a constant $D_N$ with
\begin{equation}
\label{eq:udecay}
  |u(z)|\le D_N(1+|z|)^{-N}\qquad(|\Iml z|\le\tfrac12).
\end{equation}
Choose $Y>|x_0|$ with $D_1/(1+Y)\le\tfrac12|u(z_0)|$.  Then
\begin{equation}
\label{eq:half}
  |u(z)|\le\tfrac12|u(z_0)|\qquad
  \text{whenever }|\Iml z|\le\tfrac12\text{ and }|\Rel z|\ge Y.
\end{equation}

\emph{The zeros of bounded height.}  Let $\mathcal W$ be the set of points
$z\in\mathcal Z\setminus\mathcal O_0$ with $|\Rel z|<Y$.  It is finite,
because $\mathcal W$ is bounded and $\xi$ is entire and not identically
zero, and it is invariant under
$z\mapsto-z$ and $z\mapsto\bar z$.  Put
\[
  P(z):=\prod_{w\in\mathcal W}(z-w),\qquad E(z):=(\tfrac14+z^2)\,P(z)^2 .
\]
The polynomial $P$ has real coefficients because $\mathcal W$ is invariant
under conjugation, and $P(-z)=(-1)^{|\mathcal W|}P(z)$ because $\mathcal W$
is invariant under $z\mapsto-z$.  Hence $E$ is an even polynomial with real
coefficients which vanishes on $\mathcal W$.  Moreover $E(z_0)\ne0$: one has
$z_0\notin\mathcal W$, and $\tfrac14+z_0^2=\rho_0(1-\rho_0)\ne0$ because
$\xi(0)=\xi(1)\ne0$.

\emph{The test functions.}  Since $z_0^2\notin\R$, the map
$(\alpha,\beta)\mapsto\alpha+\beta z_0^2$ is an isomorphism of real vector
spaces $\R^2\to\C$; let $\theta$ be a constant with
$|\alpha|+|\beta|\le\theta\,|\alpha+\beta z_0^2|$ for all real
$\alpha,\beta$.  For an integer $N\ge1$ let $\eta^{\star N}$ be the $N$-fold
convolution power of $\eta$, a real even function in $C_c^\infty(\R)$ with
$\widetilde{\eta^{\star N}}=u^N$ by Lemma~\ref{lem:testfn}(ii).  Choose real
numbers $\alpha_N^\pm,\beta_N^\pm$ such that
\[
  E(z_0)\,u(z_0)^N\bigl(\alpha_N^++\beta_N^+z_0^2\bigr)=i,\qquad
  iz_0\,E(z_0)\,u(z_0)^N\bigl(\alpha_N^-+\beta_N^-z_0^2\bigr)=1,
\]
and put
\[
  q_N^+(z):=E(z)\bigl(\alpha_N^++\beta_N^+z^2\bigr),\qquad
  q_N^-(z):=iz\,E(z)\bigl(\alpha_N^-+\beta_N^-z^2\bigr),\qquad
  F_{N,\pm}:=q_N^\pm(\partial)\,\eta^{\star N}.
\]
By Lemma~\ref{lem:Pi},
\begin{equation}
\label{eq:GN}
  V_N^\pm:=\widetilde{F_{N,\pm}}=q_N^\pm\,u^N,\qquad
  V_N^+(z_0)=i,\qquad V_N^-(z_0)=1 .
\end{equation}
The operator $q_N^+(\partial)$ is a polynomial with real coefficients in
$\partial^2=-d^2/dt^2$, and $q_N^-(\partial)$ is the product of
$i\partial=-d/dt$ and such a polynomial.  Therefore $F_{N,+}$ is real and
even and $F_{N,-}$ is real and odd; both lie in $C_c^\infty(\R)$, and both
are nonzero by \eqref{eq:GN}.  Since $q_N^\pm$ is divisible by
$\tfrac14+z^2$, both are primitive by Lemma~\ref{lem:Pi}.

It remains to evaluate \eqref{eq:zsum} for $F=F_{N,\pm}$.  We split
$\mathcal Z$ into $\mathcal O_0$, $\mathcal W$ and the set $\mathcal Z_Y$ of
points with $|\Rel z|\ge Y$; this is a partition because $|x_0|<Y$.

\emph{The points of $\mathcal O_0$.}  The function $V_N^+$ is even, so its
term in \eqref{eq:zsum} is $m_0V_N^+(z)^2$, and by \eqref{eq:realF}
$V_N^+(\bar z)=\overline{V_N^+(z)}$.  Hence $V_N^+(z)^2$ equals $i^2=-1$ at
$\pm z_0$ and $\overline{-1}=-1$ at $\pm\bar z_0$.  The function $V_N^-$ is
odd, so its term in \eqref{eq:zsum} is $-m_0V_N^-(z)^2$, and by
\eqref{eq:realF} $V_N^-(\bar z)=-\overline{V_N^-(z)}$.  Hence $-V_N^-(z)^2$
equals $-1$ at each of the four points.  In both cases the points of
$\mathcal O_0$ contribute $-4m_0$ to \eqref{eq:zsum}.

\emph{The points of $\mathcal W$.}  Here $E$ vanishes, so $V_N^\pm$ vanishes,
and the contribution is $0$.

\emph{The points of $\mathcal Z_Y$.}  Let $d:=\deg E+3$ and let $C_E$ be a
constant with $|E(z)|(1+|z|)^3\le C_E(1+|z|)^d$ for all $z$.  By the choice
of the coefficients and of $\theta$,
\[
  |\alpha_N^\pm|+|\beta_N^\pm|\le\theta_1\,|u(z_0)|^{-N},\qquad
  \theta_1:=\frac{\theta\max(1,|z_0|^{-1})}{|E(z_0)|},
\]
so that $|q_N^\pm(z)|\le\theta_1C_E(1+|z|)^d|u(z_0)|^{-N}$ for all $z$.  Let
$z\in\mathcal Z_Y$.  Then $|\Iml z|\le\tfrac12$, and by \eqref{eq:half}
\[
  |V_N^\pm(z)|\le\theta_1C_E(1+|z|)^d\Bigl(\frac{|u(z)|}{|u(z_0)|}\Bigr)^N
  \le\theta_1C_E\,2^{-(N-1)}(1+|z|)^d\,\frac{|u(z)|}{|u(z_0)|}.
\]
The same bound holds at $-z$, which also lies in $\mathcal Z_Y$, and
$|u(-z)|=|u(z)|$.  With \eqref{eq:udecay} for the exponent $d+1$ we obtain
\[
  \sum_{z\in\mathcal Z_Y}m(z)\,\bigl|V_N^\pm(z)\,V_N^\pm(-z)\bigr|
  \le4^{-(N-1)}\,\frac{\theta_1^2C_E^2D_{d+1}^2}{|u(z_0)|^2}
     \sum_{z\in\mathcal Z}\frac{m(z)}{(1+|z|)^2},
\]
and the last sum is finite because $\sum_\rho|\rho|^{-2}$ converges and
$|z|\ge|\rho|-\frac12$.  None of the constants on the right depends on $N$.

Adding the three contributions,
\[
  \sum_\rho\Phi_{F_{N,\pm}}(\rho)\Phi_{F_{N,\pm}}(1-\rho)
  =-4m_0+O\bigl(4^{-N}\bigr)\qquad(N\to\infty),
\]
which is negative for $N$ large.  For such an $N$ put $F_\pm:=F_{N,\pm}$.
Since $F_\pm$ is primitive, \eqref{eq:comparisonprim} gives
$\Il(\Gamma_{F_\pm},\Gamma_{F_\pm})
=-\sum_\rho\Phi_{F_\pm}(\rho)\Phi_{F_\pm}(1-\rho)>0$.
\end{proof}

We turn to the form $I^{(c)}_S$ of a fixed chart.

\begin{theorem}
\label{thm:radical}
Let $S$ be a chart and $c\in\R$.  For all $h,g\in C_c^\infty(\R)$,
\begin{equation}
\label{eq:radform}
  I^{(c)}_S(\Gamma_h,\Gamma_g)
  =\frac1{2\pi}\int_\R\hat h(\tau)\,\overline{\hat g(\tau)}\,
     \bigl(\widehat T_S(\cz,\tau)+c\bigr)\,d\tau ,
\end{equation}
the integral being absolutely convergent.  If $h\in C_c^\infty(\R)$ and
$I^{(c)}_S(\Gamma_h,\Gamma_g)=0$ for every $g\in C_c^\infty(\R)$, then $h=0$.
Thus the restriction of $I^{(c)}_S$ to the subalgebra $\DAr(S)$ of radial
classes is nondegenerate.
\end{theorem}

\begin{proof}
By \eqref{eq:radalg}, $\Gamma_h\star\Gamma_g^\dagger=\Gamma_k$ with
$k=h\star g^\dagger$, and $\hat k=\hat h\,\overline{\hat g}$.  By
\eqref{eq:plancherel} and \eqref{eq:profile},
\[
  I^{(c)}_S(\Gamma_h,\Gamma_g)=T^{(c)}_S(\Gamma_k)
  =\frac1{2\pi}\int_\R\hat k(-\tau)\bigl(\widehat T_S(\cz,\tau)+c\bigr)d\tau ,
\]
absolutely convergent, and the substitution $\tau\mapsto-\tau$ gives
\eqref{eq:radform}, because $\widehat T_S(\cz,\cdot)$ is even by
Proposition~\ref{prop:symbol}(d).

Suppose that $I^{(c)}_S(\Gamma_h,\Gamma_g)=0$ for every $g$, and put
$v(\tau):=\hat h(\tau)\bigl(\widehat T_S(\cz,\tau)+c\bigr)$.  The function
$v$ is continuous, and it is square integrable, because $\hat h$ is a
Schwartz function and $|\widehat T_S(\cz,\tau)|\le C\log(2+|\tau|)$ by
Proposition~\ref{prop:symbol}(d).  By \eqref{eq:radform} it is orthogonal in
$L^2(\R)$ to $\hat g$ for every $g\in C_c^\infty(\R)$.  These functions are
dense in $L^2(\R)$, because $C_c^\infty(\R)$ is dense in $L^2(\R)$ and the
Fourier transform is a constant multiple of a unitary operator of $L^2(\R)$,
by Plancherel's theorem.  Hence $v=0$ almost everywhere, and everywhere by
continuity.  By Proposition~\ref{prop:symbol}(b) the zeros of
$\widehat T_S(\cz,\cdot)+c$ form a discrete set.  Hence $\hat h$ vanishes off
a discrete set, so $\hat h=0$ by continuity, and $h=0$.
\end{proof}

The last assertion of Theorem~\ref{thm:radical} also follows from the
nondegeneracy of $I^{(c)}_S$ on $\DA(S)$ \cite[Theorem~6.5]{P43}, by
averaging the second argument over $K_S$.  For $G\in\DA(S)$ put
$G^\natural(w):=\int_{K_S}G(wk)\,d\mu(k)$.  This is an element of $\DA(S)$
invariant under $K_S$, so $G^\natural=\Gamma_g$ for some
$g\in C_c^\infty(\R)$ by Proposition~\ref{prop:radalg}, and
$\widehat{G^\natural}(\chi)=\widehat G(\chi)\int_{K_S}\chi\,d\mu$ is equal to
$\widehat G$ on $\cz$ and vanishes on the other components.  Since
$\widehat{\Gamma_h}$ vanishes off $\cz$, the formula
$I^{(c)}_S(F,G)=\int\widehat F\,\overline{\widehat G}\,(\widehat T_S+c)\,d\chi$
of \cite[(69)]{P43} gives
$I^{(c)}_S(\Gamma_h,G)=I^{(c)}_S(\Gamma_h,\Gamma_g)$.  Hence, if $\Gamma_h$ is
orthogonal to all radial classes, it is orthogonal to $\DA(S)$, and
$\Gamma_h=0$.

\begin{thebibliography}{9}

\bibitem{Barner1981} K.~Barner,
\emph{On A.~Weil's explicit formula},
J.~Reine Angew. Math. \textbf{323} (1981), 139--152.

\bibitem{Connes99} A.~Connes,
\emph{Trace formula in noncommutative geometry and the zeros of the Riemann
zeta function}, Selecta Math. (N.S.) \textbf{5} (1999), 29--106.

\bibitem{Davenport} H.~Davenport,
\emph{Multiplicative Number Theory}, 3rd ed., revised by
H.~L.~Montgomery, Graduate Texts in Mathematics \textbf{74}, Springer, New York, 2000.

\bibitem{NIST} F.~W.~J.~Olver, D.~W.~Lozier, R.~F.~Boisvert and C.~W.~Clark
(eds.), \emph{NIST Handbook of Mathematical Functions},
Cambridge University Press, Cambridge, 2010.

\bibitem{P43} D.~A.~Trejo Pizzo,
\emph{A correspondence algebra with a canonical trace on the idele class
group}, preprint.

\bibitem{P44} D.~A.~Trejo Pizzo,
\emph{An arithmetic Frobenius family on the idele class group}, preprint.

\bibitem{Weil1952} A.~Weil,
\emph{Sur les ``formules explicites'' de la th\'eorie des nombres premiers},
Comm. S\'em. Math. Univ. Lund (1952), Tome suppl\'ementaire, 252--265.

\end{thebibliography}

\end{document}
